\scanpage{35}
\begin{center}- 45 -\end{center}
\hfill n$^\circ$ 272

est un groupe de type fini, de rang (sur $\underline{\underline Q}$) égal à 1. Comme c'est aussi un groupe sans torsion, $\Gamma_{v'}$ est nécessairement isomorphe à $\underline{\underline Z}$, C.Q.F.D

[On pourrait aussi déduire ce corollaire directement de la prop.1, mais il me paraît plus naturel de le laisser en compagnie des cor.2 et 3].

2. - \uline{Application} : \uline{prolongement d'une place à valeurs dans un corps algébriquement clos}.

PROPOSITION 3 - \uline{Soit} $L$ \uline{une extension algébrique d'un corps}, $K$, \uline{et soit} $\theta$ \uline{une place de} $L$. \uline{L'extension} $\theta(L)/\theta(K)$ \uline{est alors une extension algébrique}.

Ce n'est que la traduction dans le langage des places de la deuxième partie de la prop.2.

De cette proposition, nous allons déduire :

PROPOSITION 4 - \uline{Soient} $K$ \uline{un corps}, $S$ \uline{un sous-anneau de} $K$, \uline{et} $g$ \uline{un homomorphisme de} $S$ \uline{dans un corps algébriquement clos} $k$. \uline{Il existe alors une place} $\theta$ \uline{de} $K$ \uline{à valeurs dans} $k$ \uline{qui prolonge} $g$.

D'après la prop.3 du \S\ 2, il existe une place $\varphi$ de $K$ à valeurs dans une extension $L$ de \inserted{$k$} qui prolonge $g$. Si $\psi$ est une place de $L$ à valeurs dans $k$, qui est l'identité sur $k$, la place composée $\theta=\psi\circ\varphi$ répond évidemment à la question. Tout revient donc à démontrer l'existence de $\psi$, autrement \unsure{doit|dit} à prouver le lemme suivant :

\uline{Lemme} 1 - \uline{Pour toute extension} $L$ \uline{d'un corps algébriquement clos} $k$, \uline{il existe une place} $\psi$ \uline{de} $L$ \uline{à valeurs dans} $k$ \uline{qui est l'identité sur} $k$.

Soit $\Omega$ l'ensemble des couples $(F,\psi)$, où $F$ est une sous-extension de $L$, et $\psi$ est une place de $F$ à valeurs dans $k$ qui est l'identité sur $k$. Si l'on ordonne \inserted{$\Omega$} par la relation de prolongement, on voit tout de suite que $\Omega$ est inductif ; comme $\Omega$ contient le couple formé de $k$ et de
\mspage{34}
